\documentclass[10pt,a4paper]{amsart}
\usepackage{amsmath,amssymb,amsthm,graphicx}
\usepackage{verbatim}
\newenvironment{reference}{\comment}{\endcomment}
%\newenvironment{reference}{}{}
\newenvironment{slogan}{\comment}{\endcomment}
\newenvironment{history}{\comment}{\endcomment}

% For commutative diagrams we use Xy-pic
\usepackage[all]{xy}

% We use 2cell for 2-commutative diagrams.
\xyoption{2cell}
% Theorem environments.
%
\theoremstyle{plain}
\newtheorem{theorem}[subsection]{Theorem}
\newtheorem{proposition}[subsection]{Proposition}
\newtheorem{lemma}[subsection]{Lemma}

\theoremstyle{definition}
\newtheorem{definition}[subsection]{Definition}
\newtheorem{example}[subsection]{Example}
\newtheorem{exercise}[subsection]{Exercise}
\newtheorem{situation}[subsection]{Situation}

\theoremstyle{remark}
\newtheorem{remark}[subsection]{Remark}
\newtheorem{remarks}[subsection]{Remarks}

\numberwithin{equation}{subsection}

% Macros
%
\def\lim{\mathop{\mathrm{lim}}\nolimits}
\def\colim{\mathop{\mathrm{colim}}\nolimits}
\def\Spec{\mathop{\mathrm{Spec}}}
\def\Hom{\mathop{\mathrm{Hom}}\nolimits}
\def\Ext{\mathop{\mathrm{Ext}}\nolimits}
\def\SheafHom{\mathop{\mathcal{H}\!\mathit{om}}\nolimits}
\def\SheafExt{\mathop{\mathcal{E}\!\mathit{xt}}\nolimits}
\def\Sch{\mathit{Sch}}
\def\Mor{\mathop{\mathrm{Mor}}\nolimits}
\def\Ob{\mathop{\mathrm{Ob}}\nolimits}
\def\Sh{\mathop{\mathit{Sh}}\nolimits}
\def\NL{\mathop{N\!L}\nolimits}
\def\CH{\mathop{\mathrm{CH}}\nolimits}
\def\proetale{{pro\text{-}\acute{e}tale}}
\def\etale{{\acute{e}tale}}
\def\QCoh{\mathit{QCoh}}
\def\Ker{\mathop{\mathrm{Ker}}}
\def\Im{\mathop{\mathrm{Im}}}
\def\Coker{\mathop{\mathrm{Coker}}}
\def\Coim{\mathop{\mathrm{Coim}}}

% Boxtimes
%
\DeclareMathSymbol{\boxtimes}{\mathbin}{AMSa}{"02}

%
% Macros for moduli stacks/spaces
%
\def\QCohstack{\mathcal{QC}\!\mathit{oh}}
\def\Cohstack{\mathcal{C}\!\mathit{oh}}
\def\Spacesstack{\mathcal{S}\!\mathit{paces}}
\def\Quotfunctor{\mathrm{Quot}}
\def\Hilbfunctor{\mathrm{Hilb}}
\def\Curvesstack{\mathcal{C}\!\mathit{urves}}
\def\Polarizedstack{\mathcal{P}\!\mathit{olarized}}
\def\Complexesstack{\mathcal{C}\!\mathit{omplexes}}
% \Pic is the operator that assigns to X its picard group, usage \Pic(X)
% \Picardstack_{X/B} denotes the Picard stack of X over B
% \Picardfunctor_{X/B} denotes the Picard functor of X over B
\def\Pic{\mathop{\mathrm{Pic}}\nolimits}
\def\Picardstack{\mathcal{P}\!\mathit{ic}}
\def\Picardfunctor{\mathrm{Pic}}
\def\Deformationcategory{\mathcal{D}\!\mathit{ef}}
\usepackage{cleveref}
\crefname{originallemma}{original lemma}{original lemmas}
\Crefname{originallemma}{Original lemma}{Original lemmas}
\makeatletter
\newlabel{algebra-lemma-yet-another-variant-local-criterion-flatness}{{0GEB}{source}{}{sourcepointer-1}{}}
\newlabel{algebra-lemma-yet-another-variant-local-criterion-flatness@cref}{{[originallemma][0][]0GEB}{[1][0][]source}{}{}{}}
\makeatother
\crefname{originalreference}{original reference}{original references}
\Crefname{originalreference}{Original reference}{Original references}
\makeatletter
\newlabel{lemma-rig-closed-point-after-localization}{{0GGG}{source}{}{sourcepointer-2}{}}
\newlabel{lemma-rig-closed-point-after-localization@cref}{{[originalreference][0][]0GGG}{[1][0][]source}{}{}{}}
\makeatother
\crefname{originalreference}{original reference}{original references}
\Crefname{originalreference}{Original reference}{Original references}
\makeatletter
\newlabel{lemma-rig-closed-point-etale}{{0GGI}{source}{}{sourcepointer-3}{}}
\newlabel{lemma-rig-closed-point-etale@cref}{{[originalreference][0][]0GGI}{[1][0][]source}{}{}{}}
\makeatother
\crefname{originalreference}{original reference}{original references}
\Crefname{originalreference}{Original reference}{Original references}
\makeatletter
\newlabel{lemma-rig-closed-point-relative}{{0GGB}{source}{}{sourcepointer-4}{}}
\newlabel{lemma-rig-closed-point-relative@cref}{{[originalreference][0][]0GGB}{[1][0][]source}{}{}{}}
\makeatother
\title{Rig-flatness is etale-local on the base}
\author{The Stacks Project Authors}
\date{stacks-project:restricted.tex:local-snapshot-20261012; GNU FDL 1.2 or later}
\begin{document}
\maketitle
\noindent\textit{Editorial note.} All delivered mathematical text is the author's own, copied from the pinned Stacks Project chapter file restricted.tex; no statement or proof step was rewritten, reordered or removed.\par\medskip
\noindent\textit{Editorial note.} The counted proof is the author's complete argument: the proof is complete: two reduction steps plus an application of the local criterion for flatness with the generating condition discharged by the rig-closed-point-etale lemma\par\medskip
\noindent\textit{Editorial note.} The label lemma-rig-flat-local-etale names the counted statement, and the proof environment immediately following it in the source is delivered verbatim as its proof.\par\medskip
\noindent\textit{Editorial note.} The counted statement is carried by the author's own lemma environment in the chapter source.\par\medskip
\noindent\textit{Allowed standard prerequisites (editorial).}\begin{itemize}
\item the WAdm\textasciicircum{}Noeth category of adic Noetherian rings with continuous adic maps topologically of finite type
\item rig-flat versus naively rig-flat via rig-closed primes
\item localization at elements f in A and the completion-of-localization description of (B')\_\{f\}
\item faithful flatness of local homomorphisms of local rings that are flat
\item the local criterion for flatness with an etale extension where the prime of the base generates the maximal ideal
\end{itemize}\medskip
\ifdefined\hypertarget\hypertarget{sourcepointer-1}{}\fi
\noindent\textit{Source reference.} Original lemma 0GEB: sources/stacks-chow/tags.txt, source lines 20471--20471. Prerequisite pointer invoked at proof line(s) 3958; it names a standard fact of a sibling chapter that the counted argument uses, and every step of the counted proof built on it is retained verbatim in the delivered body.\par\medskip
\ifdefined\hypertarget\hypertarget{sourcepointer-2}{}\fi
\noindent\textit{Source reference.} Original reference 0GGG: sources/stacks-restricted/pinned-original-source.tex, source lines 3488--3488. Prerequisite pointer invoked at proof line(s) 3922; it names a standard fact of the same chapter that the counted argument uses, and every step of the counted proof built on it is retained verbatim in the delivered body.\par\medskip
\ifdefined\hypertarget\hypertarget{sourcepointer-3}{}\fi
\noindent\textit{Source reference.} Original reference 0GGI: sources/stacks-restricted/pinned-original-source.tex, source lines 3576--3576. Prerequisite pointer invoked at proof line(s) 3971; it names a standard fact of the same chapter that the counted argument uses, and every step of the counted proof built on it is retained verbatim in the delivered body.\par\medskip
\ifdefined\hypertarget\hypertarget{sourcepointer-4}{}\fi
\noindent\textit{Source reference.} Original reference 0GGB: sources/stacks-restricted/pinned-original-source.tex, source lines 3333--3333. Prerequisite pointer invoked at proof line(s) 3956; it names a standard fact of the same chapter that the counted argument uses, and every step of the counted proof built on it is retained verbatim in the delivered body.\par\medskip
\begin{lemma}
\label{lemma-rig-flat-local-etale}
Consider a commutative diagram
$$
\xymatrix{
B \ar[r] & B' \\
A \ar[r] \ar[u]^\varphi & A' \ar[u]_{\varphi'}
}
$$
in $\textit{WAdm}^{Noeth}$ with all arrows adic and topologically
of finite type. Assume $A \to A'$ and $B \to B'$ are flat.
Let $I \subset A$ be an ideal of definition.
If $\varphi$ is rig-flat and $A/I \to A'/IA'$
is \'etale, then $\varphi'$ is rig-flat.
\end{lemma}

\begin{proof}
Given $f \in A'$ the assumptions of the lemma remain true for the digram
$$
\xymatrix{
B \ar[r] & (B')_{\{f\}} \\
A \ar[r] \ar[u]^\varphi & (A')_{\{f\}} \ar[u]
}
$$
Hence it suffices to prove that $\varphi'$ is naively rig-flat.

\medskip\noindent
Take a rig-closed prime ideal $\mathfrak q' \subset B'$.
We have to show that $(B')_{\mathfrak q'}$ is flat over $A'$.
We can choose an $f \in A$ which maps to a unit of $B'/\mathfrak q'$
such that the induced prime ideal $\mathfrak p''$ of $A_{\{f\}}$ 
is rig-closed, see Lemma \ref{lemma-rig-closed-point-after-localization}.
To be precise, here $\mathfrak q'' = \mathfrak q' B'_{\{f\}}$ and
$\mathfrak p'' = A_{\{f\}} \cap \mathfrak q''$.
Consider the diagram
$$
\xymatrix{
B'_{\mathfrak q'} \ar[r] &
(B'_{\{f\}})_{\mathfrak q''} \\
A \ar[r] \ar[u] &
A_{\{f\}} \ar[u]
}
$$
We want to show that the left vertical arrow is flat.
The top horizontal arrow is faithfully flat as it is a local
homomorphism of local rings and flat as $B'_{\{f\}}$
is the completion of a localization of the Noetherian ring
$B'_f$. Similarly the bottom horizontal arrow is flat.
Hence it suffices to prove that the right vertical arrow is flat.
Finally, all the assumptions of the lemma remain true for the diagram
$$
\xymatrix{
B_{\{f\}} \ar[r] &
B'_{\{f\}} \\
A_{\{f\}} \ar[r] \ar[u] &
A'_{\{f\}} \ar[u]
}
$$
This reduces us to the case discussed in the next paragraph.

\medskip\noindent
Take a rig-closed prime ideal $\mathfrak q' \subset B'$
and assume $\mathfrak p = A \cap \mathfrak q'$ is rig-closed as well.
This implies also the primes $\mathfrak q = B \cap \mathfrak q'$
and $\mathfrak p' = A' \cap \mathfrak q'$ are rig-closed, see
Lemma \ref{lemma-rig-closed-point-relative}.
We are going to apply
Algebra, Lemma \ref{algebra-lemma-yet-another-variant-local-criterion-flatness}
to the diagram
$$
\xymatrix{
B_\mathfrak q \ar[r] &
B'_{\mathfrak q'} \\
A_\mathfrak p \ar[u] \ar[r] &
A'_{\mathfrak p'} \ar[u]
}
$$
with $M = B_\mathfrak q$. The only assumption that hasn't been checked yet
is the fact that $\mathfrak p$ generates the maximal ideal of
$A'_{\mathfrak p'}$. This follows from
Lemma \ref{lemma-rig-closed-point-etale}.
\end{proof}

\end{document}
